\begin{pastexam}[label={px:ntust-114-ice-01}]{臺科大資工所 114 學年度 第 1 題}
（複選；每格 5 分）

(5\%) Select all the correct statements.
\begin{enumerate}[label=(\alph*)]
\item An iterative method will call itself many times using loops.
\item $10N^2+100N+10000=O(N^2)$
\item $10N^2+100N+10000=O(N^4)$
\item $10N^2+100N+10000=\Omega(N^2)$
\end{enumerate}
\end{pastexam}

\begin{pastsoln}
答案是 (b)、(c)、(d)。 令 $f(N) = 10N^2 + 100N + 10000$，漸近符號的定義見資結書第~1.5~節。
\begin{itemize}
\item (a) 錯。 迭代（iterative）的方法用迴圈重複做事，\textbf{不}呼叫自己；呼叫自己的是遞迴
（recursive）。 這句把兩者混在一起。
\item (b) 對。 $N \ge 1$ 時 $100N \le 100N^2$、$10000 \le 10000N^2$，所以
$f(N) \le (10 + 100 + 10000)N^2 = 10110N^2$，取 $c = 10110$、$N_0 = 1$。
\item (c) 對。 $O$ 是上界，$N \ge 1$ 時 $N^2 \le N^4$，由 (b) 得 $f(N) \le 10110N^4$。 上界不必是最緊的。
\item (d) 對。 $N \ge 1$ 時 $f(N) \ge 10N^2$，取 $c = 10$、$N_0 = 1$。 （(b)、(d) 合起來就是
$f(N) = \Theta(N^2)$。）
\end{itemize}
\end{pastsoln}
